\section{El contrato \texttt{.proto} y Protocol Buffers}

Protocol Buffers es un formato para definir y serializar datos estructurados
que se usa junto con gRPC~\cite{niswar2024}
para que dos procesos intercambien mensajes sabiendo de antemano qué campos
existen y de qué tipo es cada uno. A diferencia del JSON de texto que usa
nuestra implementación TCP, donde la estructura no está escrita en ningún
archivo y el servidor accede a cada clave a mano, en Protocol Buffers la
estructura y los tipos se declaran en un archivo \texttt{.proto} y los mensajes
viajan en binario. Esto no elimina la validación de los valores: nuestro
servidor gRPC sigue rechazando el emisor vacío y el \texttt{timestamp}
negativo. Otra diferencia es el tratamiento de los campos ausentes: en
\texttt{proto3} un campo que no se envía llega con su valor por defecto
(cadena vacía o 0), mientras que en TCP la clave faltante provoca un
\texttt{KeyError} que el servidor detecta y responde como error. Además, gRPC
transporta estos mensajes sobre HTTP/2~\cite{rfc9113}, que a su vez funciona
sobre TCP~\cite{rfc9293}.

\subsection{El archivo \texttt{servicio.proto}}

El archivo \texttt{proto/servicio.proto} es el contrato entre el servidor y el
cliente gRPC: describe qué mensajes existen y qué función se puede llamar de
forma remota. Su contenido, sin comentarios, es el siguiente:

\begin{verbatim}
syntax = "proto3";

package mensajeria;

service ComunicacionNodos {
  rpc EnviarMensaje(Mensaje) returns (Acuse);
}

message Mensaje {
  string sender = 1;
  int64 timestamp = 2;
  string message = 3;
}

message Acuse {
  string sender = 1;
  int64 received_timestamp = 2;
  int64 timestamp = 3;
  string message = 4;
}
\end{verbatim}

Cada parte cumple una función:

\begin{itemize}
  \item \texttt{syntax = "proto3"} indica la versión del lenguaje de
        Protocol Buffers.
  \item \texttt{package mensajeria} agrupa las definiciones y evita choques
        de nombres con otros archivos.
  \item \texttt{service ComunicacionNodos} declara la función remota
        \texttt{EnviarMensaje}: recibe un \texttt{Mensaje} y devuelve un
        \texttt{Acuse}.
  \item \texttt{Mensaje} es lo que envía el cliente y equivale al JSON
        \texttt{\{sender, timestamp, message\}} de TCP.
  \item \texttt{Acuse} es lo que responde el servidor y equivale a la
        respuesta JSON de TCP.
  \item Los números \texttt{= 1}, \texttt{= 2}, \dots{} identifican cada campo
        dentro del mensaje binario. En el binario no viajan los nombres de los
        campos, solo estos números.
\end{itemize}

\subsection{Tamaño del mensaje: Protocol Buffers frente a JSON}

Para comparar el tamaño se codificó el mismo mensaje de prueba
(\texttt{sender="nodo-A"}, \texttt{timestamp=5}, \texttt{message="hola"}) en
las dos formas. En nuestro mensaje de prueba, Protocol Buffers ocupa 16 bytes
y el JSON generado con \texttt{json.dumps} ocupa 55 bytes. Estas cifras
corresponden solo al contenido del mensaje, sin contar las cabeceras de gRPC ni
de TCP, y no se pueden generalizar a otros mensajes: con textos largos, la
diferencia relativa será menor, porque el contenido del texto ocupa lo mismo
en ambos formatos.

Otros trabajos han comparado gRPC con REST. En uno, gRPC obtuvo mejor
rendimiento en transferencias de datos grandes, pero la diferencia no resultó
estadísticamente significativa~\cite{ain2025}; en otro, las aplicaciones con
gRPC fueron, en general, más eficientes y escalables, aunque consumieron más
CPU~\cite{jcsi2024}. Sobre la serialización, se ha señalado que reducir el
tamaño del mensaje influye más en la latencia de extremo a extremo que la
velocidad de serialización~\cite{maltsev2024}. Nuestra comparación solo mide
el tamaño en bytes, no la latencia.

\subsection{Elección del tipo \texttt{int64}}

Los tres campos de tiempo guardan el valor del reloj de Lamport. Aunque este
reloj nunca es negativo y \texttt{uint64} también sería válido, se eligió
\texttt{int64} por dos razones. Primero, facilita la compatibilidad con
lenguajes sin enteros sin signo nativos, como Java, si en el futuro se escribe
un cliente en otro lenguaje. Segundo, con \texttt{int64} el servidor puede
recibir un valor negativo y rechazarlo de forma controlada con
\texttt{INVALID\_ARGUMENT}, como comprueba la prueba
\texttt{test\_timestamp\_negativo}. En Python ambos tipos se leen como
\texttt{int}, así que la elección no cambia el resto del código del servidor.

\subsection{Generación del código}

El archivo \texttt{.proto} no se ejecuta. Se procesa con
\texttt{grpc\_tools.protoc} desde la raíz del repositorio:

\begin{verbatim}
python -m grpc_tools.protoc -I. \
    --python_out=. --grpc_python_out=. \
    proto/servicio.proto
\end{verbatim}

Este comando genera dos archivos:

\begin{itemize}
  \item \texttt{servicio\_pb2.py}: las clases \texttt{Mensaje} y
        \texttt{Acuse} para crear y leer mensajes.
  \item \texttt{servicio\_pb2\_grpc.py}: la clase base
        \texttt{ComunicacionNodosServicer}, de la que hereda nuestro servidor,
        y el \texttt{ComunicacionNodosStub}, que usa el cliente para llamar a
        \texttt{EnviarMensaje} como si fuera una función local.
\end{itemize}

La opción \texttt{-I.} toma la raíz del repositorio como punto de partida, de
modo que el import generado queda como \texttt{from proto import servicio\_pb2}.